% 第 6 章：有限动作认证恢复
\section{有限动作认证恢复}\label{sec:res-certified}

\subsection{功能定位}
认证恢复（\srcpath{include/hacdcpf/resilience/certified\_restoration.hpp}）把恢复表述为
\textbf{有限动作目录}上的主问题 MIP：每个候选动作先过信息物理/MESS 可执行性检查，再由
多保真动态 DAE 证书验证，最后以 no-good 割迭代。协调器
\srcpath{CertifiedRestorationCoordinator::solve} 消费动作集
\fld{CertifiedRestorationAction} 返回 \fld{CertifiedRestorationResult}，
\fld{model\_scope}=\texttt{"finite-action restoration master MIP + cyber/MESS executability + dynamic DAE oracle"}。

\subsection{动作与可执行性}
\fld{CertifiedRestorationAction} 携带目标值、动态事件序列
（\srcpath{dynamics::DynamicEvent}）、信息物理命令链契约（\fld{command\_path\_available}/
\fld{authorization\_valid}/\fld{acknowledgement\_available}）与 MESS 契约（\fld{mess\_travel\_time\_s}、
\fld{mess\_connection\_deadline\_s}、\fld{mess\_required\_energy\_mwh}/\fld{mess\_available\_energy\_mwh}、
\fld{mess\_grid\_forming\_capable}）。
\texttt{MultiFidelity\allowbreak Certificate\allowbreak Engine::checkExecutability}
返回 \fld{ExecutabilityReport}（\fld{cyber\_executable}$\wedge$\fld{mess\_executable}$\Rightarrow$\fld{executable}）。

\subsection{证书量：残差与雅可比}
对轨迹快照，原生相域残差评估器给出微分与代数残差
\begin{equation}
r_f=\dot x-f(t,x,y),\qquad
r_g=\begin{bmatrix}\Re(I_{ac}-Y_{ac}V_{abc})\\ \Im(I_{ac}-Y_{ac}V_{abc})\\ I_{dc}-G_{dc}V_{dc}\end{bmatrix}.
\end{equation}
以 $\widetilde x=D_x^{-1}x$ 缩放，稠密有限差分雅可比构成 Schur 约简局部矩阵
\begin{equation}
A_D=D_x^{-1}\big(F_x-F_yG_y^{-1}G_x\big)D_x.
\end{equation}

\subsection{重构缺陷与输出定向界}
L1/L2 用同一物理 DAE、不同积分步：每对接受端点构造三次 Hermite 重构 $\widehat x_H(t)$，
其连续缺陷与累积量
\begin{equation}
\rho_H(t)=D_x^{-1}\big[\dot{\widehat x}_H(t)-f(t,\widehat x_H,\psi(\widehat x_H))\big],\quad
R_x=\kappa_g e_0+\int\|\rho_H\|_\infty\,dt+\kappa_g\int\|r_g\|_\infty\,dt.
\end{equation}
全态指示器仅供审计：$\eta=K_\Phi(T)R_x$，$K_\Phi(T)=\sup_{0\le\tau\le T}\|\exp(A_D\tau)\|_\infty$。
分类用输出定向界：对安全输出 $q_j$，
\begin{equation}
c_j^\mathsf{T}=(q_{j,x}-q_{j,y}G_y^{-1}G_x)D_x,\quad
\Delta_j(t)=K_j(t,0)e_0+\int_0^tK_j(t,s)\|\rho_H(s)\|_\infty ds+B_j\|r_g\|_{\infty,[0,t]},
\end{equation}
$K_j(t,s)=\|c_j^\mathsf{T}(t)\Phi_D(t,s)\|_1$、$B_j=\|q_{j,y}G_y^{-1}\|_1$。COI 频率的直接代数项为 0，
AC 电压保留。逐动作证书
\texttt{MultiFidelity\allowbreak Certificate} 报告 \fld{margins}（含 \fld{Delta\_j}）、
重构缺陷积分、$\eta$ 与构造方法。

\subsection{保真级与阈值}
\begin{itemize}
  \item \textbf{L1}：完整器件模型 + 粗积分步（诊断）；
  \item \textbf{L2}：完整器件模型 + 细积分步（诊断）；
  \item \textbf{L3}：完整 DAE 步 + 情景/时域阈值 oracle——频率
        $\ge$\fld{min\_frequency\_hz}（49）、RoCoF $\le$\fld{max\_rocof\_hz\_s}（1）、
        AC 电压 $\ge$\fld{min\_ac\_voltage\_pu}（0.90），并含换流器电流观测门。
\end{itemize}
\fld{proof\_valid} 仅表示 L3 分类输入完整，\textbf{不}意味全局解析稳定证明；L1/L2 常量为采样
估计，\fld{constants\_are\_global\_bounds}=false，即使区间分隔零也只是诊断。

\subsection{主问题迭代与目录条件最优}
HiGHS 主问题每次选恰好一个可执行、未排除动作；仅当 MIP 间隙为零时用作上界；L3-safe 动作
更新在位解；仅对 \fld{proof\_valid}=true 且标签 \texttt{unsafe} 的 L3 结果加 no-good 割；
unresolved/failed 停止循环而不加割。若每次主解间隙为零且每个 L3 标签在其声明情景/时域上
可靠，归纳可得：返回动作在\textbf{所编码有限目录}内最优（\fld{optimality\_proven\_over\_catalog}），
但不覆盖目录外动作或 L3 省略的物理。结果记 \fld{master\_mip\_solves}、\fld{dynamic\_oracle\_calls}、
\fld{no\_good\_cuts} 与逐轮 \fld{iterations}（\fld{MasterIterationRecord}）。
